% Laget automatisk av elevark/lag_elevark.py fra programmering/python-variable-krefter.html. Endringer her blir overskrevet.
\documentclass[11pt]{article}
\usepackage{elevark}
\setlength{\svarhoyde}{1.4cm}
\begin{document}
\elevarkhode{Fysikk 1 · Krefter}{Variable krefter i Python}{https://lektorbo.no/programmering/python-variable-krefter.html}{Koden på nett}

Bevegelseslikningene gjelder bare når akselerasjonen er konstant. Hvis kraften endrer seg underveis, for eksempel fordi friksjonen eller luftmotstanden endrer seg, kan vi i stedet regne ut bevegelsen i mange små tidssteg.

\section*{Oppgaver}

\begin{oppgaveliste}
\oppgave \fase{test}{Feilsøk}Kjør programmet i del 1. Hva skjer? Finn feilen og rett den.
\svar
\oppgave \fase{test}{Regn}Hva blir bremselengden og bremsetiden? Sammenlign med en bil som bremser med konstant $μ = 0{,}20$. Regn ut den for hånd med bevegelseslikningene, og sjekk med programmet.
\svar
\oppgave \fase{forklar}{Forklar}Hvorfor kan vi ikke bruke bevegelseslikningene i del 1? Hvorfor skriver vi $a = -R(s)/m$ inne i løkka og ikke før den?
\svar
\oppgave \fase{test}{Kjør}Kjør del 2. Hvor høyt kommer ballen, og hvor fort faller den til slutt? Sammenlign med et kast uten luftmotstand $(k = 0)$.
\svar
\oppgave \fase{forklar}{Forklar}Hvorfor er akselerasjonen større enn $g$ på vei opp og mindre enn $g$ på vei ned? Hva er akselerasjonen i toppunktet?
\svar
\oppgave \fase{forklar}{Forklar}Hvorfor bruker vi $v \cdot |v|$ og ikke $v^{2}$ i luftmotstanden?
\svar
\oppgave \fase{utfordring}{Utfordring}Hvor stor må $k$ være for at ballen bare skal komme 1,0~m over kastepunktet? Prøv deg fram, eller lag en løkke som øker $k$ i små steg.
\svar
\end{oppgaveliste}

\end{document}
